\section{Local estimates for the determinant}
\label{s4-local}

This section follows Section~4 of the source, pp.~9--13, with
$D_N^6$, $h=K-N$, and the scalar $S_K$ unchanged.  Write $v_p^G$ for
the minimum valuation of the coefficients of a polynomial in $X$.
The object estimated here is $\Delta_{9,K}$ before multiplication by
$S_K$.  Fix $M\geq40$.  The inner and outer bases serve different
purposes: the first spreads vanishing among residue classes, while the
second cancels the poles whose harmonic sums contain multiples of $p$.

\subsection{The inner range and its new row weights}
\label{s4-inner}

Assume, as in source (4.1), that
\begin{equation}
 K\in40\mathbb Z_{>0},\qquad K\geq200M^2,
 \qquad K/M<p\leq K/3.
 \label{s4-inner-range}
\end{equation}
Then $p>200M$ and $p^2>5K$.  Put
\[
 m_A=\lfloor A/p\rfloor,\qquad
 \ell_A(a)=\#\{1\leq j\leq A:j\equiv\pm a\pmod p\},
 \qquad 1\leq a\leq m=(p-1)/2.
\]
If a row polynomial vanishes to order $\nu_i(a)$ at $t=-a^2$,
the ordinary-class bound replacing source (4.2) is
\begin{equation}
 \nu_i(a)+\nu_j(a)+6\ell_N(a)-\ell_K(a)-8.
 \label{s4-ordinary-source}
\end{equation}
At zero the bound replacing (4.3) is unchanged:
\begin{equation}
 2\nu_i(0)+2\nu_j(0)+12m_N-2m_K+1.
 \label{s4-zero-source}
\end{equation}
\textbf{[Derived]}
Indeed, the distribution factor is now $p^{-8}$.  At an ordinary
source $x=a+pz$, $x^9$ is a unit to leading order, giving the loss
$-8$.  At $x=pz$, its factor $p^9$ compensates for that distribution
factor and leaves $+1$.
\textbf{[Formal analogue]} Applying these counts to the entire rational
summand, including its polynomial quotient, requires the analytic
completion estimate from Section~3.  We keep that transplant explicit.

Retain the source allocation, namely
\begin{equation}
 \begin{gathered}
 L_0=4M+10,\qquad b_a=3\ell_N(a),\\
 mT+E=h-L_0+3(N-m_N),\qquad 0\leq E<m,\\
 L_a=T-b_a+\epsilon_a,\qquad Z_a=L_a+b_a=T+\epsilon_a.
 \end{gathered}
 \label{s4-allocation}
\end{equation}
Here $\epsilon_a=1$ for the first $E$ classes in decreasing order of
$\ell_K(a)$, and is zero otherwise.  With $x=K/p$, the unchanged
inequalities
$2Hx-21/20<T<2Hx$ and $b_a\leq6\alpha x+3$ give
$T-b_a>2\lambda x-81/20\geq3/2$.  Thus the dimensions are
nonnegative and sum to $h$.

Use the polynomials from source (4.5):
\begin{equation}
 E_{a,i}(t)=\prod_{\substack{0\leq c\leq m\\c\ne a}}
 (t+c^2)^{L_c}(t+a^2)^i,
 \qquad 0\leq i<L_a.
 \label{s4-inner-basis}
\end{equation}
Their degrees are less than $h$.  Distinct factors are coprime modulo
$p$, and their local Taylor-coordinate matrices are triangular with
unit diagonals.  Hence this is a $\mathbb Z_p$-unimodular basis.
The weights replacing source (4.6)--(4.8) are
\begin{align}
 w^{(9)}_{a,i}&=i+b_a-\frac{\ell_K(a)+8}{2},\qquad a>0,
 \label{s4-inner-weights}\\
 w^{(9)}_{0,i}&=\min\left\{
 2i+6m_N-m_K+\frac12,
 \min_{1\leq c\leq m}\left(Z_c-\frac{\ell_K(c)+8}{2}\right)
 \right\},
 \label{s4-inner-zero-weight}\\
 \gamma^{\mathrm{in}}_{9,p}&=
 2\sum_{a=0}^{m}\sum_{i=0}^{L_a-1}w^{(9)}_{a,i}.
 \label{s4-inner-gamma}
\end{align}

The comparison between two ordinary blocks reduces, just as before,
to
$Z_c-Z_a+1+(\ell_K(a)-\ell_K(c))/2\geq0$.
The counts take two consecutive values, and the ordering of the
extras resolves the only potentially adverse case.  At zero the
half-weight of an ordinary row is at least $7M+41/2$, comfortably
above its assigned weight.  Formula~\eqref{s4-inner-zero-weight}
handles the zero-block rows at every source.

The local degree checks also survive.  At an ordinary source the
leading numerator degree is at most $2(T+1)<23M/5+2$; at zero it is
at most
\[
 9+4L_0+12m_N\leq169M/10+49<p+1.
\]
\textbf{[Conditional]}
Subject to the Section~3 analytic completion estimate, every matrix
entry in this basis has valuation at least the sum of its row
weights.  Expanding the determinant therefore gives
\begin{equation}
 v_p^G(\Delta_{9,K})\geq\gamma^{\mathrm{in}}_{9,p}.
 \label{s4-inner-bound}
\end{equation}
Relative to the source weights, every ordinary weight decreases by
two; every zero-block weight decreases by an amount between zero
and two.  Consequently the useful exact comparison is
\begin{equation}
 -4h\leq\gamma^{\mathrm{in}}_{9,p}-\gamma^{\mathrm{in}}_{5,p}
 \leq-4h+4L_0.
 \label{s4-inner-shift}
\end{equation}
Thus the limiting inner cost numerator changes by $4\lambda x$,
with only an $O_M(1)$ error in the finite exponent.

\subsection{Removing the outer polynomial correction}
\label{s4-rank-correction}

Use the source outer hypotheses, with the harmless strengthened
threshold $p\geq11$:
\begin{equation}
 p\leq K<3p,\qquad p^2>2K,\qquad 2N<p,
 \qquad5N\leq2p-2.
 \label{s4-outer-range}
\end{equation}
\textbf{[Derived]}
The denominator $40320$ is a $p$-adic unit.  Von Staudt--Clausen
therefore gives $v_p(\mu_9(t^e))\geq-1$.
If $e<4p-5$ and a Bernoulli denominator contains $p$, then
$2e+2=j(p-1)$ for some $1\leq j\leq7$.
The factor $2e+2+j=jp$ occurs among $2e+3,\ldots,2e+9$ and
cancels that denominator.  Hence
\begin{equation}
 \mu_9(t^e)\in\mathbb Z_p\qquad(e<4p-5).
 \label{s4-moment-cutoff}
\end{equation}

Choose $c_e\in\mathbb Z_p$ representing $p\mu_9(t^e)$ modulo $p$,
with $c_e=0$ below this cutoff.  Replace the polynomial moments by
$\mu_9^0(t^e)=\mu_9(t^e)-c_e/p$, leaving simple-pole values
unchanged.  Set $W=D_N^5$ and $D_{\mathrm{tail}}=D_K/D_N$;
the rational input is still $Wt^{i+j}/D_{\mathrm{tail}}$.
Its polynomial quotient has degree at most $i+j+6N-K$.  Therefore
\begin{equation}
 G_{9,K}=A+p^{-1}L,\qquad
 \operatorname{rank}_{\mathbb Q_p}L\leq
 r_{9,p}:=\max(0,K+4N-4p+4),
 \label{s4-matrix-rank}
\end{equation}
where $L$ is integral.  Indeed,
$L_{ij}=0$ if $i+j<K-6N+4p-5$, so its first $h-r_{9,p}$
rows and columns vanish exactly.  This is a rank bound over
$\mathbb Q_p$, rather than only after reduction modulo $p$.
For our fixed parameters and $p>K/3$,
$K+4N-4p+4<4-K/30$; thus $r_{9,p}=0$ throughout the outer
range once $K\geq120$.

\subsection{The outer basis, class costs, and rank lemma}
\label{s4-outer-basis}

For an ordinary square class put $\ell=\ell_K(a)$ and
$\delta=\mathbf1_{a\leq N}$.  Let $Q_a$ contain its remaining
$\ell-\delta$ factors, $P_a=D_{\mathrm{tail}}/Q_a$, and let
$E_a$ contain its $\ell-2$ factors with $j>p$.  The rows are
$P_aq_{a,i}$, with
\begin{equation}
 q_{a,i}(t)=
 \begin{cases}
 (t+a^2)^i,&0\leq i<\ell-2,\\
 E_a(t)(t+a^2)^{i-(\ell-2)},&\ell-2\leq i<\ell-\delta.
 \end{cases}
 \label{s4-outer-polynomials}
\end{equation}
Together with the zero-class rows these form a unimodular basis:
each local sequence is monic of successive degrees, and distinct
classes have unit resultants.  Cross-class entries of $A$ are
integral because their quotients are integral polynomials.

\textbf{[Derived]}
For two low indices, the row factors supply $i+j$, $W$ supplies
$5\delta$, the node derivative loses $\ell-\delta-1$, and the
harmonic value loses at most nine powers of $p$.  The resulting
bound is $i+j+6\delta-\ell-8$.  A high-index row cancels all
poles with $j>p$.  The remaining values at $a,p-a$ are integral
and congruent modulo $p$: use
$H_{p-a}^{(9)}\equiv H_{a-1}^{(9)}\pmod p$ and the displayed
simple-pole formula for $\mu_{9,X}$.  The $X$ coefficients agree
because $(p-a)^8\equiv a^8$.  The constant-coefficient difference
also vanishes: the extra $a^{-9}$ in $H_a^{(9)}$ cancels the
change in $1/(2a)$.  The divided difference is consequently
integral, including after multiplication by polynomial numerators.

Thus the analogue of source (4.12) is
\begin{equation}
 w^{(9)}_{a,i}=
 \begin{cases}
 \min\{0,i+3\delta-(\ell+8)/2\},&i<\ell-2,\\
 0,&i\geq\ell-2.
 \end{cases}
 \label{s4-outer-weights}
\end{equation}
At zero a single pole has weight $-1/2$; for $p,2p$, rows
$1,t+p^2$ give entry bounds $-3,-1,0$, so weights $-2,0$
remain valid.  Here each pole value has valuation at least $-1$
because $j^8H_j^{(9)}$ loses only one power when $p\mid j$.
The spacing of the two squared nodes costs two further powers.

For reference, the elementary rank lemma from source (4.13)
is unchanged.  If $v(A_{ij})\geq w_i+w_j$, the $w_i$ are
nonpositive half-integers, $L$ is integral of rank at most $r$,
and exactly $z$ weights are zero, then
\begin{equation}
 v\det(A+p^{-1}L)\geq2\sum_iw_i-\min(r,z).
 \label{s4-rank-lemma}
\end{equation}
A complementary-minor term using $k$ rows and columns of $L$
has valuation at least
$2\sum_iw_i-k-\sum_{i\in I}w_i-\sum_{j\in J}w_j$.
Terms with $k>r$ vanish.  Ordering weights decreasingly, the
largest possible loss is
$\max_{0\leq k\leq\min(r,h)}(k+2\sum_{i=1}^kw_{(i)})$:
each zero weight increases it by one, and every subsequent
increment is nonpositive.  This proves the lemma, also for
Gauss valuations.  Unimodular changes of basis preserve the
integrality and rank of $L$.

\subsection{Explicit outer exponent and the primes beyond \texorpdfstring{$K$}{K}}
\label{s4-outer-formula}

\textbf{[Derived]}
For a removed class, the new table is
\[
\begin{array}{c|rrrrr}
 \ell&2&3&4&5&6\\\hline
 -2\sum_iw^{(9)}_{a,i}&0&5&10&15&20\\
 \#\{i:w^{(9)}_{a,i}=0\}&1&1&1&1&1
\end{array}
\]
An unremoved class costs $11(\ell-2)$ and has two zero weights.
Put
\[
 q=\lfloor K/p\rfloor\in\{1,2\},\quad v=K-qp,\quad
 u=\max(0,N+v-p+1),\quad t_p=\min(N,v)+u.
\]
Over all ordinary classes,
$\sum(\ell-2)=K-p-q+1$.  Over the $N$ removed classes it is
$2(q-1)N+t_p$: the two indicators $a\leq v$ and $a\geq p-v$
account for $t_p$.  Each removed class saves $6(\ell-2)$ from
the unremoved cost.  Adding the zero-class costs $1,4$ yields
\begin{equation}
 \gamma^{\mathrm{out}}_{9,p}=
 \begin{cases}
 -11(K-p)+6t_p-1-\min(r_{9,p},p-1-N),&K<2p,\\
 -11(K-p)+7+12N+6t_p-\min(r_{9,p},p-N),&K\geq2p.
 \end{cases}
 \label{s4-outer-gamma}
\end{equation}
The zero counts are $p-1-N$ and $p-N$, respectively.  Unlike
the source table, the removed classes never acquire additional
zero weights as $\ell$ increases.  Applying
\eqref{s4-rank-lemma} gives
$v_p^G(\Delta_{9,K})\geq\gamma^{\mathrm{out}}_{9,p}$.
As a finite consistency check, $K=40,N=3,p=29$ gives $-104$,
matching the independently computed constant-coefficient
valuation; this check does not establish the general estimate.

For $p>K$, every pole is below $p$ and each square class has at
most two poles.  The same congruence and divided-difference
argument makes its contribution integral.  The largest possible
polynomial quotient degree is $K+4N-2<4p-5$, so no polynomial
correction remains.  Consequently
\begin{equation}
 v_p^G(\Delta_{9,K})\geq0\qquad(p>K).
 \label{s4-large-prime-integrality}
\end{equation}
These local statements supply normalization exponents.  They do
not by themselves supply the final decay inequality.
